\section{Applications and Future Work}
\label{sec:future}

The construction pipeline behind OenoBench is intended to be reusable.
We outline three threads.

\subsection{Closing the wine-knowledge gap in low-cost models}
Small open-weight models trail proprietary frontier ones by 20--30\,pp
on this benchmark, while the broad-benchmark gap is only 10--15\,pp
(Section~\ref{sec:eval-cost}). The 38{,}104-fact corpus is a natural
supervised-fine-tuning target via a two-stage recipe: instruction
tuning on held-out question/source-fact pairs, then LoRA adaptation
on the full fact-question-source triples in the release schema. Both
stages run on consumer GPUs without proprietary weights or
pre-training compute. The contextual slice
(Section~\ref{sec:eval-cb}) is the harder-to-game evaluation surface,
so measured gain tracks genuine wine reasoning. A
Llama-3.3-70B-class adaptation could plausibly close most of the
GPT-5-mini gap at $\approx$1\% of proprietary inference cost.

\subsection{Case-based evaluation with industry partners}
Multiple-choice tests have ceilings: a perfect OenoBench score does
not establish useful behaviour in real workflows. We are scoping an
\textbf{applied case-based dataset} with industry partners along three
personas:
\begin{itemize}\itemsep -1pt
  \item \emph{Sommelier $\rightarrow$ customer service:} pairing
    prompts rated on recommendation, justification, and substitutions.
  \item \emph{Winemaker $\rightarrow$ blending decisions:} vintage
    write-ups, lab analyses, and regulatory constraints in context;
    rated on technical correctness, compliance, and stylistic intent.
  \item \emph{Viticulturist $\rightarrow$ vineyard decisions:}
    weather, soil/disease, canopy/yield records; recommendations
    rated against decisions actually taken in the same vintage.
\end{itemize}
The OenoBench audit-and-review tooling (Appendix~\ref{app:screenshots})
carries over directly.

\subsection{Generalising the scrape-to-audit pipeline}
The pipeline modules --- tier-of-authority taxonomy, atomic-fact
extraction, multi-strategy multi-model generation with per-model
quotas, $\kappa$-calibrated multi-agent audit, and closed-book
pre-screen --- are domain-agnostic. Three near-term targets share
wine's structural features: \emph{horticulture/agronomy} (regulated
appellations, peer-reviewed journals, certification ladders);
\emph{regulated medicine} (clinical guidelines, drug formularies;
current benchmarks rely on board-exam questions misaligned with
practice); and \emph{financial regulation} (FASB, IFRS, jurisdictional
filings; CFA/CPA ladders). The released code reduces months of
expert authoring to weeks of curation.
